\section{Baselines}  \label{sec:baselines}

We run several baselines to benchmark the performance of state-of-the-art LLM and VLM agents. All models are prompt-based and provided with 3 in-context examples (one from each environment), which share no overlap with the benchmark tasks. The prompts we use are provided in the appendix. We summarize the results in Tab.~\ref{tab:results} and describe the baselines in detail in the following sections.

\subsection{Text-based LLM Agents}

Several prior works developed autonomous agents by prompting text-based LLMs~\citep{zhou2023webarena,kim2023language,liu2023bolaa}. We benchmark several text-based LLM agents with Chain-of-Thought prompting~\cite{wei2022chain} over the accessibility tree representations of the websites as input, and leave more advanced prompting strategies for future work. We test API-based LLMs, including GPT-4 Turbo (\texttt{gpt-4-1106-preview}), GPT-3.5 Turbo (\texttt{gpt-3.5-turbo-1106}), and Gemini-Pro, as well as open sourced LLMs (LLaMA-2-70B, Mixtral-8x7B).


\begin{table*}[t]
\centering
\resizebox{\linewidth}{!}{%
\begin{tabular}{@{}lccccccc@{}}
\toprule
\multirow{2}{*}{\textbf{Model Type}} & \multirow{2}{*}{\textbf{LLM Backbone}} & \multirow{2}{*}{\textbf{Visual Backbone}} & \multirow{2}{*}{\textbf{Inputs}} & \multicolumn{4}{c}{\textbf{Success Rate ($\uparrow$)}} \\
\cmidrule(lr){5-8}
& & & & \textbf{Classifieds} & \textbf{Reddit} & \textbf{Shopping} & \textbf{Overall} \\
\midrule
\multirow{5}{*}{Text-only} & LLaMA-2-70B & \multirow{5}{*}{-} & \multirow{5}{*}{Acc. Tree} & 0.43\% & 1.43\% & 1.29\% & 1.10\% \\
                        & Mixtral-8x7B &  &  & 1.71\% & 2.86\% & 1.29\% & 1.76\% \\
                        & Gemini-Pro &  &  & 0.85\% & 0.95\% & 3.43\% &  2.20\% \\
                        & GPT-3.5 &  &  & 0.43\% & 0.95\% & 3.65\% &  2.20\% \\
                        & GPT-4 &  &  & 5.56\% & 4.76\% & 9.23\% & 7.25\% \\
\midrule
\multirow{6}{*}{Caption-augmented} & LLaMA-2-70B & BLIP-2-T5XL & \multirow{6}{*}{Acc. Tree + Caps} & 0.00\%  &  0.95\%  & 0.86\% & 0.66\% \\
                        & Mixtral-8x7B & BLIP-2-T5XL &  & 1.28\% & 0.48\% & 2.79\% & 1.87\% \\
                        & GPT-3.5 & LLaVA-7B &  & 1.28\% & 1.43\% & 4.08\% & 2.75\% \\
                        & GPT-3.5 & BLIP-2-T5XL &  & 0.85\%  & 1.43\% & 4.72\% & 2.97\% \\
                        & Gemini-Pro & BLIP-2-T5XL &  & 1.71\%  & 1.43\% & 6.01\% & 3.85\% \\
                        & GPT-4 & BLIP-2-T5XL &   & 8.55\% &  8.57\% &  16.74\% & 12.75\% \\
\midrule
\multirow{4}{*}{Multimodal} & \multicolumn{2}{c}{IDEFICS-80B-Instruct} & \multirow{4}{*}{Image + Caps + Acc. Tree} &  0.43\% & 0.95\% & 0.86\% & 0.77\% \\ 
                        & \multicolumn{2}{c}{CogVLM} &  & 0.00\% & 0.48\% & 0.43\% &  0.33\% \\ 
                        % & \multicolumn{2}{c}{Qwen-VL} &  & \color{blue}{XX\%} &  \color{blue}{XX\%} & \color{blue}{XX\%}  &  \color{blue}{XX\%} \\ 
                        & \multicolumn{2}{c}{Gemini-Pro} &  & 3.42\% & 4.29\% & 8.15\% &  6.04\% \\ 
                        & \multicolumn{2}{c}{GPT-4V} &  & 8.12\%   &  12.38\% &  \textbf{19.74\%} & 15.05\% \\
\midrule
\multirow{4}{*}{Multimodal (SoM)} &  \multicolumn{2}{c}{IDEFICS-80B-Instruct} & \multirow{4}{*}{Image + Caps + SoM} & 0.85\% & 0.95\% & 1.07\% & 0.99\% \\ 
                        & \multicolumn{2}{c}{CogVLM} &  & 0.00\% & 0.48\% & 0.43\% &  0.33\% \\ 
                        % & \multicolumn{2}{c}{Qwen-VL} &  & \color{blue}{XX\%} & \color{blue}{XX\%} & \color{blue}{XX\%} &  \color{blue}{XX\%} \\ 
                        & \multicolumn{2}{c}{Gemini-Pro} &  & 3.42\% & 3.81\% & 7.73\% & 5.71\% \\ 
                        % & \multicolumn{2}{c}{Gemini-Flash-1.5} &  & 3.85\% & 4.76\% & 8.80\% &  6.59\% \\ 
                        % & \multicolumn{2}{c}{Gemini-Pro-1.5} &  & 5.98\% & 12.86\% & 14.59\% &  11.98\% \\ 
                        & \multicolumn{2}{c}{GPT-4V} &  &  \textbf{9.83\%} &  \textbf{17.14\%} & 19.31\% & \textbf{16.37\%}  \\ 
                        % & \multicolumn{2}{c}{GPT-4o} &  &  \textbf{17.95\%} &  \textbf{17.14\%} & \textbf{21.24\%} & \textbf{19.45\%}  \\ 
\midrule
\multirow{1}{*}{Human Performance} &   - & - & \multirow{1}{*}{Webpage} & 91.07\% &  87.10\% & 88.39\% &  88.70\% \\ 
\bottomrule
\end{tabular}
}
% \vspace{-0.12in}
\caption{Success rates of baseline LLM and VLM agents on \BenchmarkName{}.}  \label{tab:results}
% \vspace{-0.15in}
\end{table*}


\subsection{Image Caption Augmented LLM Agents}

\BenchmarkName{} is a visually grounded benchmark, and we expect that leveraging complementary visual information would improve performance. Hence, we run pretrained image captioning models on every \texttt{img} element on the HTML page, and augment the accessibility tree with this information as the image alt-text before passing this as input to the LLM agents. If a task contains input images, we also caption them and include the captions as part of the prompt. We run experiments on GPT-3.5 with two recent image captioning models, BLIP-2-T5XL~\citep{li2023blip} and LLaVA-v1.5-7B~\citep{liu2023improved}. Our results with GPT-3.5 as the LLM backbone (``Caption-augmented'' section of Tab.~\ref{tab:results}) suggest that the LLaVA and BLIP-2 captioning models achieve comparable performance. Since BLIP-2 achieves a slightly higher success rate, is a smaller model, and requires less GPU VRAM, we use it as the captioning backbone for the remaining experiments.

\subsection{Multimodal Agents} \label{sec:som_method}


Finally, we benchmark strong API-based and open-source VLMs as agents. 
% Unlike the previous text-based models, these multimodal models are trained on large datasets of paired text and images, allowing them to learn joint representations between vision and language. 
We evaluate several models capable of processing multiple interleaved image-and-text inputs: GPT-4V~\citep{openai2023gpt4}, Gemini-Pro~\citep{team2023gemini}, \href{https://huggingface.co/HuggingFaceM4/idefics-80b-instruct}{IDEFICS-80B-Instruct} (a reimplementation of Flamingo~\citep{alayrac2022flamingo}), and CogVLM~\citep{wang2023cogvlm}. 
We experiment with two settings:

\paragraph{Image Screenshot + Captions + Accessibility Tree:} This approach provides the accessibility tree representation augmented with image captions as accessibility tree alt-text from BLIP-2-T5XL (similar to the caption-augmented agent), as well as the screenshot of the current webpage as inputs. This provides the model with both the structural information and the visual context of the website.

\paragraph{Image Screenshot + Captions + SoM:} Inspired by Set-of-Marks prompting~\cite{yang2023som}, we perform an initial preprocessing step by using JavaScript to automatically annotate every interactable element on the webpage with a bounding box and a unique ID. The annotated screenshot containing bounding boxes and IDs, are provided as input to the multimodal model along with a text representation of the SoM (see Fig.~\ref{fig:som_figure}). Similar to the baselines above, we also provide the captions from BLIP-2-T5XL for all \texttt{img} elements on the page. There have been several projects\footnote{\href{https://github.com/ddupont808/GPT-4V-Act}{GPT-4V-ACT} and \href{https://github.com/ishan0102/vimGPT}{vimGPT} propose similar interfaces.} that propose similar representations. Most have been proof-of-concept demos, and to the best of our knowledge, we are the first to systematically benchmark this on a realistic and interactive web environment. % such as \BenchmarkName{}.

